% Section 6, part 6.
\subsection{The lattice theorem}
\label{sec:mod-lattice}

\begin{lem}[Counting]
\label{lem:mod-count}
For every $w\geq 0$, $V_w$ has exactly $\mathrm d(w)$ basis triples $(m,k,b)$ with $b\neq 1$ and exactly $\mathrm e(w)$ with $b=1$.
\end{lem}

\begin{proof}
Fix $m$ with $5m\leq w$ and put $r:=w-5m$. A triple with $b\neq 1$ needs $w_b+4k=r$ with $w_b\in\{1,4,3,2\}$ for $b=0,2,3,4$. These four weights meet every class modulo $4$ exactly once, so there is exactly one candidate $b$, and it works unless $w_b>r$; this happens only for the class $0$, where $w_2=4$ and $r=0$, that is $w=5m$. Hence the number is $\lfloor w/5\rfloor+1-[5\mid w]=\mathrm d(w)$. The triples with $b=1$ ($w_1=0$) are the pairs $(m,k)$ with $5m+4k=w$.
\end{proof}

\begin{lem}[The leg class]
\label{lem:mod-classL}
\begin{enumerate}
\item $\bbQ[p,\bar p]=\bbQ[p,x]=\bbQ[v][p]+\bbQ[v][x]$, and $\bbQ[p,x]$ is a free $\bbQ[v,x]$-module with basis $1,p,p^2,p^3,p^4$.
\item $\bbQ[v][p^{\pm1}]\cap\bbQ(v)[x]=\bbQ[v][x]$.
\item $\cL\cap\bbQ[v][p^{\pm1}]=\bbQ[p,\bar p]$: a Laurent polynomial with coefficients in $\bbQ[v]$ lies in $\cL$ if and only if each of its terms $v^mp^j$ has $j\geq -m$.
\item $\cL/\bbQ(v)[x]=\bbQ[p,\bar p]/\bbQ[v,x]=\bbQ[v][p]/\bbQ[v]=\bigoplus_{i=1}^4\bbQ[v,x]\,p^i$; it is free over $\bbQ[v]$, of dimension $\mathrm d(w)$ in weight $w$.
\end{enumerate}
\end{lem}

\begin{proof}
(i) $\bar p=x-p^4/4$ gives $\bbQ[p,\bar p]=\bbQ[p,x]$, and $v=p\bar p=px-p^5/4$. The elements $p$ and $x$ are algebraically independent, since $p$ and $v$ are and $v=px-p^5/4$. The map $\bbQ[v,x][T]/(T^5-4xT+4v)\to\bbQ[p,x]$, $T\mapsto p$, is onto; its source is $\bbQ[x,T]$ (eliminate $v=(4xT-T^5)/4$), so by the independence it is an isomorphism, which gives the basis. Next, $\bbQ[v][p]+\bbQ[v][x]\subset\bbQ[p,\bar p]$. Conversely let $F\in\bbQ[p,\bar p]$ be homogeneous with a negative power of $p$, and let $v^mp^{-n}$ ($n>0$, $m\geq n$) be its term with the most negative exponent; by homogeneity it is the only term with this exponent. Since $x^n=\sum_t\binom nt4^{t-n}v^tp^{4n-5t}$ has the lowest term $v^np^{-n}$, subtracting a multiple of $v^{m-n}x^n$ gives an element of $\bbQ[p,\bar p]$ whose most negative exponent is larger; by induction $F\in\bbQ[v][p]+\bbQ[v][x]$. (ii) Let $P=\sum_{n\leq N}a_nx^n$ with $a_N\neq 0$ lie in $\bbQ[v][p^{\pm1}]$. Only $x^N$ contributes to $p^{4N}$, with the coefficient $a_N4^{-N}$, so $a_N\in\bbQ[v]$; induct on $N$. (iii) The inclusion $\supset$ follows from (i). If $F=R+P$ with $R\in\bbQ[v][p]$ and $P\in\bbQ(v)[x]$ is a Laurent polynomial over $\bbQ[v]$, then $P=F-R\in\bbQ[v][x]$ by (ii), so $F\in\bbQ[p,\bar p]$ by (i). (iv) The map $\bbQ[p,\bar p]\to\cL/\bbQ(v)[x]$ is onto because $\bbQ[v][p]\subset\bbQ[p,\bar p]$, with kernel $\bbQ[p,\bar p]\cap\bbQ(v)[x]=\bbQ[v,x]$ by (ii). By (i), $\bbQ[p,\bar p]/\bbQ[v,x]=\bigoplus_{i=1}^4\bbQ[v,x]p^i$. As the lowest term of $x^n$ is $v^np^{-n}$, $\bbQ[v][p]\cap\bbQ[v][x]=\bbQ[v]$, so the quotient is also $\bbQ[v][p]/\bbQ[v]$, whose weight-$w$ part has the basis $v^mp^{w-5m}$, $0\leq m\leq w/5$, $5m<w$: $\mathrm d(w)$ elements.
\end{proof}

\begin{lem}[Reduction modulo $v$]
\label{lem:mod-modv}
In every weight, $\Lambda\otimes\bbQ[v]/(v)\colon\mathbb F/v\mathbb F\to\mathsf N/v\mathsf N$ is an isomorphism; here $\Lambda$ is followed by the projection $\mathsf M\to\mathsf N$.
\end{lem}

\begin{proof}
$\Lambda$ is graded and $\bbQ[v]$-linear with values in $\mathsf M$ (\cref{lem:mod-closure}). \emph{Source.} In weight $w$, $\mathbb F/v\mathbb F$ has the basis $\psi^kH^b$ with $w_b+4k=w$: one element if $4\nmid w$; $\{\psi^{n-1}H^2,\psi^nH\}$ if $w=4n\geq 4$; $\{H\}$ if $w=0$. \emph{Target.} $\mathsf N/v\mathsf N=\bbQ[p,\bar p]/(\bbQ[v,x]+p\bar p\,\bbQ[p,\bar p])\oplus\log p\cdot\bbQ[x]$. The ring $\bbQ[p,\bar p]$ is the polynomial ring in $p,\bar p$, and $\bbQ[p,\bar p]/(p\bar p)$ has the basis $1$, $p^a$, $\bar p^c$ ($a,c\geq 1$). The image of $\bbQ[v,x]$ is spanned by $1$ and the $x^n\equiv p^{4n}/4^n+\bar p^n$; so the first summand has the basis $\{p^a:a\geq 1\}$, with $\bar p^n\equiv -4^{-n}p^{4n}$, and $(\mathsf N/v\mathsf N)_w$ has the basis $\{p^w\}$ if $w\geq 1$, $4\nmid w$; $\{p^{4n},\log p\cdot x^n\}$ if $w=4n\geq 4$; $\{\log p\}$ if $w=0$. \emph{The map.} By \cref{lem:mod-closure}~(iii), $I$ induces a map on $\mathsf N/v\mathsf N$. By \cref{lem:mod-closure}~(i), $I(p^a)\equiv -p^{a+4}/(a+4)$ modulo $v\mathsf M$ for $a\geq 1$, so $\phi^b_k\equiv(-1)^kp^{w_b+4k}/\prod_{i=1}^k(w_b+4i)$ for $b=0,4,3$; since $\bar p\equiv -p^4/4$ modulo $\bbQ[v,x]$, $\phi^2_k\equiv(-1)^{k+1}p^{4k+4}/\big(4\prod_{i=1}^k(4+4i)\big)$; and by \cref{lem:mod-closure}~(ii) the $\log p$-component of $\phi^1_k$ is $(-1)^kx^k/k!$ modulo $v$. In weight $4n$ the matrix in these bases is triangular with non-zero diagonal entries; in the other weights it is a non-zero $1\times 1$ matrix (the factors $\varrho_b^{-1}$ do not matter).
\end{proof}

\begin{thm}[The lattice theorem]
\label{thm:mod-lattice}
For every $w\geq 0$:
\begin{enumerate}
\item[(a)] $V_w$ has $\mathrm d(w)$ basis triples with $b\neq 1$ and $\mathrm e(w)$ with $b=1$;
\item[(b)] for $X\in\sKo\cap V_w$, $\Lambda(X)$ lies in $\bbQ[p,\bar p]=\bbQ[v][p]+\bbQ[v][x]$, exactly and not only modulo functions of $x$;
\item[(c)] $\Lambda\colon\sKo\cap V_w\to(\cL/\bbQ(v)[x])_w$ is bijective, and both sides have dimension $\mathrm d(w)$;
\item[(d)] $\Lambda\colon V_w\to\mathsf N_w$ is bijective; $\mathsf N$ is a free $\bbQ[v]$-module with basis $\{\Lambda(\psi^kH^b)\}$.
\end{enumerate}
The coefficients must lie in $\bbQ[v]$: over $\bbQ(v)$ the statement fails, since $\Lambda(H^2/v)=-1/p\notin\cL$.
\end{thm}

\begin{proof}
(a) is \cref{lem:mod-count}. (b) $\Lambda(X)\in\mathsf M$ by \cref{lem:mod-closure}, and it is log-free by \cref{prop:mod-dii}, so $\Lambda(X)\in\bbQ[p,\bar p]$, which equals $\bbQ[v][p]+\bbQ[v][x]$ by \cref{lem:mod-classL}~(i). (d) $\Lambda\colon\mathbb F\to\mathsf N$ is a graded $\bbQ[v]$-linear map between modules with finite-dimensional homogeneous pieces, zero in negative weights. Let $C$ be its cokernel. By right exactness of $\otimes\bbQ[v]/(v)$ and \cref{lem:mod-modv}, $C/vC=0$. If $C\neq 0$, let $w$ be minimal with $C_w\neq 0$; then $(vC)_w=vC_{w-5}=0$, so $C_w=(C/vC)_w=0$, a contradiction. So $\Lambda$ is onto. Since $\dim V_w=\mathrm d(w)+\mathrm e(w)$ (\cref{lem:mod-count}) and $\dim\mathsf N_w=\mathrm d(w)+\mathrm e(w)$ (\cref{lem:mod-classL}~(iv), and $\dim\bbQ[v,x]_w=\mathrm e(w)$ as $v,x$ are algebraically independent), $\Lambda_w$ is bijective and $\mathsf N\cong\mathbb F$ is free. (c) By (b) the image lies in $\bbQ[p,\bar p]$, and $\cL/\bbQ(v)[x]=\bbQ[p,\bar p]/\bbQ[v,x]$ by \cref{lem:mod-classL}~(iv). If $\Lambda(X)\in\bbQ(v)[x]$, then $\Lambda(X)\in\bbQ[p,\bar p]\cap\bbQ(v)[x]=\bbQ[v,x]$ (\cref{lem:mod-classL}~(ii)), so $\Lambda(X)=0$ in $\mathsf N$ and $X=0$ by (d). Given $\phi\in(\bbQ[p,\bar p]/\bbQ[v,x])_w$, the element $X:=\Lambda^{-1}(\phi\oplus 0)$ of (d) has a log-free leg, so $X\in\sKo$ by \cref{prop:mod-dii}. Both sides have dimension $\mathrm d(w)$ by \cref{lem:mod-classL}~(iv). Finally $\Lambda(H^2)=\varrho_2^{-1}v/p=-v/p$, and $-1/p\notin\cL$ by \cref{lem:mod-classL}~(iii).
\end{proof}

\begin{cor}[Splitting of $V_w$]
\label{cor:mod-split}
$V_w=(\sKo\cap V_w)\oplus\operatorname{span}\{v^m\psi^kH:5m+4k=w\}$. Every $X\in V_w$ is $X_0+\sum a_{m,k}v^m\psi^kH$ with $X_0\in\sKo$, the $a_{m,k}$ being read off from the $\log p$-coefficient of $\Lambda(X)$.
\end{cor}

\begin{proof}
The decomposition is \cref{lem:tra-frames}~(a) in weight $w$. Alternatively: by \cref{lem:mod-closure}~(ii) and induction, the $\log p$-coefficient of $\phi^1_k$ is $(-1)^kx^k/k!$ plus terms of lower degree in $x$, so the map that sends $X$ to the $\log p$-coefficient of $\Lambda(X)$, restricted to the span of the $v^m\psi^kH$, is triangular onto $\bbQ[v,x]_w$ ($v^m\psi^kH\leftrightarrow v^mx^k$), hence bijective; its kernel on $V_w$ is $\sKo\cap V_w$ by \cref{prop:mod-dii}.
\end{proof}

\begin{cor}[Splitting of $\sK\cap V_w$]
\label{cor:mod-Ksplit}
$\sK\cap V_w=(\sKo\cap V_w)\oplus\bbQ v^{w/5}H$ if $5\mid w$, and $\sK\cap V_w=\sKo\cap V_w$ otherwise. On $\sK$ the only legs beyond $\bbQ[p,\bar p]$ are multiples of $\log p$: for $X\in\sK\cap V_w$, $\Lambda(X)=\Lambda(X-\Pi_0(X)H)+\Pi_0(X)\log p$ with $\Lambda(X-\Pi_0(X)H)\in\bbQ[p,\bar p]$.
\end{cor}

\begin{proof}
By \cref{prop:mod-dii}, $X\in\sK$ if and only if the $\log p$-coefficient $\sum_i\Pi_i(X)(-x)^i/i!$ of $\Lambda(X)$ does not involve $x$, that is, equals $\Pi_0(X)$, which lies in $\bbQ v^{w/5}$ (\cref{def:tra-tower}; zero unless $5\mid w$). Write $X=X_0+Y$ as in \cref{cor:mod-split}. Then $\Lambda(Y)$ has the same $\log p$-coefficient as $\Lambda(X)$, while $\Lambda(\Pi_0(X)H)=\Pi_0(X)\log p$ has the $\log p$-coefficient $\Pi_0(X)$; by injectivity on the span of the $v^m\psi^kH$, $Y=\Pi_0(X)H$. Conversely $\sKo\subset\sK$ and $H\in\sK$. The last statement follows from \cref{thm:mod-lattice}~(b), as $X-\Pi_0(X)H\in\sKo\cap V_w$.
\end{proof}

\begin{cor}[A basis of $\sKo\cap V_w$]
\label{cor:mod-basis}
The elements $v^m\fr(k,a)$ ($a\neq 1$) of weight $w$ form a basis of $\sKo\cap V_w$, and $\dim\sKo\cap V_w=\mathrm d(w)$.
\end{cor}

\begin{proof}
By \cref{lem:tra-frames}~(b), (c) in the graded setting, the frame coordinates identify $\sKo\cap V_w$ with the families $(c_{(k,a)})_{a\neq 1}$ in which $c_{(k,a)}$ is a rational multiple of $v^{(w-4k-\nu_a-1)/5}$, zero unless the exponent is a non-negative integer; so the $v^m\fr(k,a)$ of weight $w$, that is the triples $(m,k,a)$ with $a\neq 1$ and $w_a+4k+5m=w$, form a basis. Their number is $\mathrm d(w)$ by \cref{thm:mod-lattice}~(a). (The same dimension follows from \cref{thm:mod-lattice}~(c).)
\end{proof}

% Section 6, part 7.
\subsection{L-regularity and the Gap Theorem}
\label{sec:mod-Lreg}

We use the closed formula of \cite[Thm.~7.1]{ABDKS25}: for an arbitrary rational function $x(z)$ on $\bbP^1$ such that $\dd x$ has only simple zeros, the differentials of topological recursion for the curve $(x=x(z),y=z)$ with the Bergman kernel and the initial form $\omega^{(0)}_{1,0}=-y\,\dd x$ are
\begin{equation}
\label{eq:mod-abdks}
\omega^{(g)}_{m,0}=[\hbar^{2g}]\sum_{r_1,\dots,r_m\geq 0}\prod_{i=1}^m\big(\mathcal D_{x_i}\big)^{r_i}[w_i^{r_i}]\sum_{\Gamma}\hbar^{2g(\Gamma)}\prod_{\{i,j\}\in E(\Gamma)}W_{i,j}\prod_{i=1}^mW_i
\end{equation}
\cite[(7.9)]{ABDKS25}, where $\Gamma$ runs over the connected simple graphs on the vertices $1,\dots,m$, $g(\Gamma)$ is the first Betti number, $\mathcal D_{x_i}\eta:=\dd(\eta/\dd x_i)$ on $1$-forms in $z_i$, $S(u):=\sinh(u/2)/(u/2)$,
\begin{gather*}
W_i:=-\exp\!\big(-w_i(S(w_i\hbar\partial_{z_i})-1)\,x(z_i)\big)\frac{\dd z_i}{w_i},\\
W_{i,j}:=\frac{1}{\hbar^2}\Big(\exp\!\big(\hbar^2w_iw_jS(w_i\hbar\partial_{z_i})S(w_j\hbar\partial_{z_j})(z_i-z_j)^{-2}\big)-1\Big)
\end{gather*}
\cite[(7.6), (7.7)]{ABDKS25}, and $W_{i,j}=\sum_{k\geq 0}(\hbar^2/4)^k(w_i+w_j)^{2k}w_iw_j(z_i-z_j)^{-2k-2}$; for one point \cite[(7.10)]{ABDKS25},
\begin{equation}
\label{eq:mod-710}
\sum_{g\geq 1}\hbar^{2g}\omega^{(g)}_{1,0}=\sum_{r\geq 0}\mathcal D_x^r[w^r]W_1 .
\end{equation}
The standing assumptions of \cite[\S1.1]{ABDKS25} (a compact curve, $x,y$ meromorphic, the critical points of $x$ simple, $y$ regular with $\dd y\neq 0$ there) hold for $x=\Wv$, $y=p$, $v\neq 0$: $\Wv$ is rational, $\dd\Wv=(p^5-v)\dd p/p^2$ has exactly five zeros, all simple, and $y=p$ has $\dd y=\dd p\neq 0$. The kernel of \cite[(1.1)]{ABDKS25} with $\omega^{(0)}_{1,0}=-y\,\dd x$ is the negative of ours (\cref{conv:mod-kernel}), so $\omega_{h,n}=(-1)^n\omega^{(h)}_{n,0}$.

\begin{lem}[Structure of the correlators]
\label{lem:mod-abdks}
For $2h-2+n>0$ there are an integer $K\geq 0$ and $N_{h,n}\in\bbQ[v][p_1,\dots,p_n]$ with
\[
\omega_{h,n}=N_{h,n}(p;v)\,\frac{\prod_i\dd p_i}{\prod_i(p_i^5-v)^K}.
\]
\end{lem}

\begin{proof}
(1) \emph{Each term of \eqref{eq:mod-abdks} lies in $\bbQ[v][p_i,p_i^{-1},(p_i^5-v)^{-1},(p_i-p_j)^{-1}]\prod_i\dd p_i$, and for fixed $(h,n)$ there are finitely many terms.} Since $S(u)-1=\sum_{k\geq 1}u^{2k}/(4^k(2k+1)!)$,
\[
-w\big(S(w\hbar\partial_p)-1\big)\Wv(p)=\sum_{k\geq 1}w^{2k+1}\hbar^{2k}\epsilon_k(p),\qquad \epsilon_k:=-\frac{\Wv^{(2k)}}{4^k(2k+1)!}\quad(\text{local letters}),
\]
and $\Wv^{(2k)}=(p^4/4)^{(2k)}+(2k)!\,v\,p^{-2k-1}\in\bbQ[v][p,p^{-1}]$. In $W_i$ every $\hbar^{2k}$ comes with $w_i^{2k+1}$, and in $W_{i,j}$ every $\hbar^2$ comes with $(w_i+w_j)^2$ and one more factor $(p_i-p_j)^{-2}$ ($z_i=p_i$). So at fixed order $\hbar^{2h}$ the degrees in the $w_i$ are bounded, only finitely many $r_i$ and graphs contribute, and the coefficients are polynomials over $\bbQ$ in the $\epsilon_k(p_i)$ and the $(p_i-p_j)^{-1}$. Finally $\mathcal D_x(f\,\dd p)=\dd\big(p^2f/(p^5-v)\big)$ keeps this ring and only adds powers of $(p^5-v)^{-1}$. (2) \emph{Clearing.} Write the finite sum as $\omega_{h,n}=P\prod_i\dd p_i\big/\big(\prod_ip_i^{a}\prod_i(p_i^5-v)^{K}\prod_{i<j}(p_i-p_j)^{e}\big)$ with $P\in\bbQ[v][p]$. By \cref{conv:mod-kernel}, $\omega_{h,n}$ has poles only at the branch points; $p_i=0$ and the diagonals contain none ($v\neq 0$). So for every $v\neq 0$ the polynomial $P$ vanishes to order at least $e$ along $p_i=p_j$ and at least $a$ along $p_i=0$. Divide $P$ by $(p_i-p_j)^e$, respectively $p_i^a$, as polynomials in $p_i$ over $\bbQ[v][p_k:k\neq i]$ (the divisors are monic, so the division algorithm works over this ring): the remainder vanishes for every $v\neq 0$ and every $p$, hence it is $0$. One factor at a time, all these denominators cancel in $\bbQ[v][p]$, and no negative power of $v$ arises.
\end{proof}

\begin{thm}[L-regularity]
\label{thm:mod-Lreg}
Let $2h-2+n>0$. For legs $\phi_1,\dots,\phi_n\in\bbQ[v][p]$ the residue pairing $\langle\phi_1\cdots\phi_n\rangle^{\mathrm{TR}}_h$ lies in $\bbQ[v]$. For legs in $\cL=\bbQ[v][p]+\bbQ(v)[x]$ the residue pairing lies in $\bbQ[v]$; the components in $\bbQ(v)[x]$ pair to zero.
\end{thm}

\begin{proof}
(1) Fix all variables but $p_l$ off the branch points. By \cref{lem:mod-abdks}, $\phi_l(p_l)\omega_{h,n}$ is a rational $1$-form in $p_l$ with poles at the branch points and possibly at $p_l=\infty$; it is holomorphic at $p_l=0$. By the residue theorem, $\sum_b\Res_{p_l=\varsigma_b}=-\Res_{p_l=\infty}$. (2) The result is again a polynomial over $\prod_{m\neq l}(p_m^5-v)^K$: in one variable, $(p^5-v)^{-K}=\sum_{j\geq 0}\binom{K+j-1}{j}v^jp^{-5K-5j}$ for $|p|^5>|v|$, so
\[
\Res_{p=\infty}\frac{p^a\,\dd p}{(p^5-v)^K}=\begin{cases}-\binom{K+j-1}{j}v^j&\text{if }a+1=5(K+j)\text{ with an integer }j\geq 0,\\ 0&\text{otherwise,}\end{cases}
\]
an element of $\bbQ[v]$. The iterated residues commute, as there are no diagonal poles. So $\langle\phi_1\cdots\phi_n\rangle^{\mathrm{TR}}_h=(-1)^n\Res_{p_1=\infty}\cdots\Res_{p_n=\infty}\big[\prod_l\phi_l(p_l)\,N_{h,n}\prod_l\dd p_l/\prod_l(p_l^5-v)^K\big]$ is a finite $\bbQ[v]$-combination of such residues, hence in $\bbQ[v]$. (3) By \cref{lem:mod-xfun} the components in $\bbQ(v)[x]$ of legs in $\cL$ pair to zero, whatever decomposition is chosen.
\end{proof}

By \cref{lem:mod-scale}~(ii) a pairing with monomial legs $p^{m_l}$ is $Cv^e$ with $e=2-2h-n+\sum_lm_l/5$; \cref{thm:mod-Lreg} says that $C=0$ whenever $e<0$.

\begin{prop}[The period of $\omega_{h,1}$]
\label{prop:mod-period}
Let $h\geq 1$. Then:
\begin{enumerate}
\item $\omega_{h,1}=\dd\Phi_h$ with $\Phi_h$ rational and holomorphic at $p=0$ and $p=\infty$;
\item $\int_\infty^0\omega_{h,1}=\frac{B_{2h}}{2h}\,v^{1-2h}$ along any path avoiding the branch points;
\item for $h\geq 2$, $\FTR_h=\dfrac{B_{2h}}{2h(2h-2)}\,v^{2-2h}$.
\end{enumerate}
\end{prop}

\begin{proof}
\emph{Step 1 (exactness).} Write $\omega^{\mathrm A}_{h,1}:=\omega^{(h)}_{1,0}=-\omega_{h,1}$. By \eqref{eq:mod-710}, $\omega^{\mathrm A}_{h,1}=[\hbar^{2h}]\sum_{r\geq 1}\mathcal D_x^r[w^r]W_1$; the term $r=0$ is $-[w^1]e^{E}\dd p=0$, where $E:=\sum_kw^{2k+1}\hbar^{2k}\epsilon_k=O(w^3)$. For $r\geq 1$, $\mathcal D_x^r(a\,\dd p)=\dd\big(D_x^{r-1}(a/x')\big)$ with $D_x=(1/x')\,\dd/\dd p$ on functions (\cref{def:mod-xi}). Hence $\omega^{\mathrm A}_{h,1}=\dd\Phi^{\mathrm A}_h$ with $\Phi^{\mathrm A}_h:=\sum_{r\geq 1}D_x^{r-1}(a_r/x')$, $a_r:=-[\hbar^{2h}][w^{r+1}]e^{E}=-\sum\frac1{j!}\prod_{i=1}^j\epsilon_{k_i}$, the sum over $j\geq 1$ and $k_1+\dots+k_j=h$ with $r=2h-1+j$; it is finite ($1\leq j\leq h$). Put $\Phi_h:=-\Phi^{\mathrm A}_h$.

\emph{Step 2 ($p=\infty$).} $\deg\epsilon_k\leq 4-2k$, $1/x'=p^2/(p^5-v)$ has degree $-3$, and $D_x$ lowers the degree by $4$; every term of $\Phi^{\mathrm A}_h$ has degree at most $(4j-2h)-3-4(2h-2+j)=5-10h\leq -5$. So $\Phi^{\mathrm A}_h$ is holomorphic at $\infty$ and $\Phi^{\mathrm A}_h(\infty)=0$.

\emph{Step 3 ($p=0$).} Write $\epsilon_k=\epsilon^0_k+(\text{terms of order}\geq 0)$ with $\epsilon^0_k:=-v\,p^{-2k-1}/(4^k(2k+1))$; further $1/x'=-(p^2/v)(1-p^5/v)^{-1}$, and $D_x=\delta_0+(\text{an operator raising the order by at least }6)$ with $\delta_0:=-(p^2/v)\,\dd/\dd p$. In the \emph{pure part} of a term, all $\epsilon_k$, $1/x'$, $D_x$ are replaced by $\epsilon^0_k$, $-p^2/v$, $\delta_0$. In the variable $\bar p=v/p$ one has $\delta_0=\dd/\dd\bar p$, $-p^2/v=-v/\bar p^2$ and $\epsilon^0_k=-\bar p^{2k+1}v^{-2k}/(4^k(2k+1))$. With the local letters $\tilde w:=w\bar p$, $\hbar':=\hbar/(2v)$, $E_0:=-\sum_{k\geq 1}\hbar'^{2k}\tilde w^{2k+1}/(2k+1)$ and $C_{h,r}:=[\hbar'^{2h}][\tilde w^{r+1}]\exp(E_0)$, the pure part of $a_r$ is $-(2v)^{-2h}\bar p^{r+1}C_{h,r}$, and the pure part of the $r$-th term is $(\dd/\dd\bar p)^{r-1}\big[(2v)^{-2h}v\,C_{h,r}\bar p^{r-1}\big]=v(2v)^{-2h}(r-1)!\,C_{h,r}$, a constant. Replacing a pure factor by its correction raises the order in $p$ by at least $3$ (by at least $2k+1$ for $\epsilon_k$, at least $5$ for $1/x'$ and for $D_x$), and $\delta_0$ never lowers the order. So every term is a constant plus $O(p^3)$: $\Phi^{\mathrm A}_h$ is holomorphic at $0$, and $\Phi^{\mathrm A}_h(0)=v(2v)^{-2h}\Sigma_h$ with $\Sigma_h:=\sum_r(r-1)!\,C_{h,r}$.

\emph{Step 4 (the identity $\Sigma_h=-4^hB_{2h}/(2h)$).} Let $\mathsf a(\mathsf x):=\sum_{k\geq 1}\mathsf x^k/(2k+1)$, so that $E_0=-\tilde w\,\mathsf a(\hbar'^2\tilde w^2)$. With $r+1=2h+j$, the coefficient $C_{h,r}$ comes from $E_0^j/j!$ alone, and $C_{h,r}=\frac{(-1)^j}{j!}[\mathsf x^h]\mathsf a(\mathsf x)^j$. Hence, by the binomial series $(1+\mathsf a)^{1-2h}=\sum_j(-1)^j\binom{2h-2+j}{j}\mathsf a^j$ (the term $j=0$ has no $\mathsf x^h$, as $h\geq 1$, and the terms $j>h$ none either),
\[
\Sigma_h=\sum_{j=1}^h(-1)^j\frac{(2h-2+j)!}{j!}[\mathsf x^h]\mathsf a^j=(2h-2)!\,[\mathsf x^h](1+\mathsf a(\mathsf x))^{1-2h}.
\]
With $\mathsf x=\mathsf y^2$, $1+\mathsf a(\mathsf y^2)=\operatorname{artanh}(\mathsf y)/\mathsf y$, so $[\mathsf x^h](1+\mathsf a)^{1-2h}=[\mathsf y^{1}]\operatorname{artanh}(\mathsf y)^{1-2h}$. By the Lagrange--B\"urmann formula for $\mathsf y=\tanh\mathsf u$, $[\mathsf y^n]H(\mathsf u(\mathsf y))=\frac1n\Res_{\mathsf u=0}H'(\mathsf u)\tanh(\mathsf u)^{-n}\dd\mathsf u$ for $n\neq 0$ and every Laurent series $H$ (change of variables and one integration by parts), applied to $H(\mathsf u)=\mathsf u^{1-2h}$ and $n=1$: $[\mathsf y^1]\mathsf u^{1-2h}=(1-2h)[\mathsf u^{2h}](\mathsf u\coth\mathsf u)=(1-2h)4^hB_{2h}/(2h)!$. So $\Sigma_h=-4^hB_{2h}/(2h)$ and $\Phi^{\mathrm A}_h(0)=-\frac{B_{2h}}{2h}v^{1-2h}$. This proves (i).

\emph{Step 5.} (ii) Since $\Phi_h$ is rational, $\int_\infty^0\omega_{h,1}=\Phi_h(0)-\Phi_h(\infty)=-\Phi^{\mathrm A}_h(0)=\frac{B_{2h}}{2h}v^{1-2h}$ along any path avoiding the branch points. (iii) $\Phi_{y\dd x}=p^5/5-v\log p$, so $(2-2h)\FTR_h=\frac15\langle p^5\rangle^{\mathrm{TR}}_h+v\langle-\log p\rangle^{\mathrm{TR}}_h$, the residue pairings with any branch of $\log p$ (\cref{def:mod-pair}). Here $\langle p^5\rangle^{\mathrm{TR}}_h\in\bbQ[v]$ by \cref{thm:mod-Lreg}, of exponent $2-2h<0$ (\cref{lem:mod-scale}~(ii)), hence $0$; and $\langle-\log p\rangle^{\mathrm{TR}}_h=\int_{\mathfrak r}\omega_{h,1}=-\int_\infty^0\omega_{h,1}=-\frac{B_{2h}}{2h}v^{1-2h}$ by (ii). So $\FTR_h=\frac{-v}{2-2h}\cdot\frac{B_{2h}}{2h}v^{1-2h}=\frac{B_{2h}}{2h(2h-2)}v^{2-2h}$.
\end{proof}

\begin{lem}[The vanishing period]
\label{lem:mod-vanper}
For $h\geq 1$, $\langle-\log p\rangle^{\mathrm{TR}}_h=\int_{\mathfrak r}\omega_{h,1}=\sum_b\Res_{\varsigma_b}(-\log p)\,\omega_{h,1}=-\frac{B_{2h}}{2h}\,v^{1-2h}$.
\end{lem}

\begin{proof}
The first equality is \cref{def:mod-pair}. For the second, $\log p$ is single-valued near each $\varsigma_b$ and $\Res_{\varsigma_b}(\log p\,\dd\Phi_h)=-\Res_{\varsigma_b}\Phi_h\,\dd p/p$; the residue theorem for the rational form $\Phi_h\,\dd p/p$, whose other poles are $0$ and $\infty$, gives $\sum_b\Res_{\varsigma_b}\log p\cdot\omega_{h,1}=\Phi_h(0)-\Phi_h(\infty)=\int_\infty^0\omega_{h,1}$, while $\int_{\mathfrak r}\omega_{h,1}=-\int_\infty^0\omega_{h,1}$ for the ray $\mathfrak r$ from $0$ to $\infty$. The value is \cref{prop:mod-period}~(ii).
\end{proof}

\begin{thm}[The Gap Theorem]
\label{thm:mod-gap}
On the model curve, for every $h\geq 2$, $\FTR_h=\dfrac{B_{2h}}{2h(2h-2)}v^{2-2h}$ in the convention of \cref{conv:mod-kernel}; hence $\Fm_h=c_hv^{2-2h}$ with $c_h=(-5)^{1-h}\dfrac{B_{2h}}{2h(2h-2)}=\dfrac{(-1)^{h-1}B_{2h}}{2h(2h-2)\,5^{h-1}}$.
\end{thm}

\begin{proof}
\Cref{prop:mod-period}~(iii) and \cref{cor:mod-trF}.
\end{proof}

\begin{rmk}[A first proof through a formula of Hock]
\label{rmk:mod-hock}
The value of $\FTR_h$ can also be read off from \cite[Cor.~1.2, (1.3)]{Hock25}, whose kernel and free energy \cite[(2.1), (1.1)]{Hock25} are those of \cref{conv:mod-kernel}; for $y$ unramified (here $y=p$ is a global coordinate) it expresses $\FTR_h$ as $\frac{1}{2-2h}\sum\Res\,[\hbar^{2h-1}]$ of an explicit series, the residues being taken at the logarithmic singularities of $x$ (there are none here) and at the other singularities of $x$ and $y$, here the poles $p=0$ and $p=\infty$. The residue at $\infty$ vanishes by a degree count, and at $p=0$ only the leading terms contribute, giving $v^{2-2h}A_h$ with $A_h=(2h-2)!\,[\mathsf x^h](1+\mathsf a(\mathsf x/4))^{1-2h}=4^{-h}\Sigma_h=-B_{2h}/(2h)$ by the Lagrange inversion of Step~4 above, so that $\FTR_h=A_hv^{2-2h}/(2-2h)$. This argument is only sketched here and is not used. \draftnote{The closed forms printed in \cite[\S3.2]{Hock25} are not used; only Cor.~1.2 there is cited.}
\end{rmk}

% Section 6, part 8.
\subsection{Theorem V: logarithmic legs}
\label{sec:mod-V}

In this subsection the model curves are considered as a family in $v\neq 0$, and $\partial_v|_x$ denotes the derivative in $v$ at fixed $x$, computed locally away from the branch points through the inverse branches $p(x,v)$; $\partial_v|_p$ is the derivative at fixed $p$. Pairings in mixed form are as in \cref{def:mod-pair}.

\begin{lem}[Fixed $p$ versus fixed $x$]
\label{lem:mod-fixedx}
$\partial_vp|_x=-p/(p^5-v)$. For every family of meromorphic forms, $\partial_v|_x=\partial_v|_p-\mathrm{Lie}_{\cX}$ in each variable, where $\cX:=\frac{p}{p^5-v}\frac{\dd}{\dd p}$. Moreover $\partial_v(y\,\dd x)|_x=-\dd p/p=\int_{q\in\mathfrak r}\omega_{0,2}(\cdot,q)$.
\end{lem}

\begin{proof}
At fixed $x$, $\partial_vp|_x=-(\partial_vx|_p)/(\partial_px)=-(1/p)\big/\big((p^5-v)/p^2\big)=-p/(p^5-v)$. For $\omega=f(p,v)\,\dd p$ the chain rule gives $\partial_v|_x\omega=(\partial_vf+f_p\,\partial_vp|_x)\dd p+f\,\dd(\partial_vp|_x)=\partial_v|_p\omega-\mathrm{Lie}_{\cX}\omega$, and the same holds in each variable of a multidifferential. Since $y=p$, $\partial_v(y\,\dd x)|_x=(\partial_vp|_x)\,\dd x=-\frac{p}{p^5-v}\cdot\frac{p^5-v}{p^2}\dd p=-\dd p/p$, and $\int_{q\in\mathfrak r}\frac{\dd p\,\dd q}{(p-q)^2}=\dd p\Big[\frac{1}{p-q}\Big]_{q=0}^{q=\infty}=-\frac{\dd p}{p}$.
\end{proof}

\begin{lem}[Integration by parts]
\label{lem:mod-ibp}
Let $w$ be a meromorphic $1$-form on $\bbP^1$ with poles only at the branch points, zero residues, holomorphic at $p=0$ and $O(\dd p/p^2)$ at $p=\infty$ (for instance $\omega_{h,n}$ in one slot, $2h-2+n>0$).
\begin{enumerate}
\item For $\phi\in\bbQ[v][p]$ let $(\cX\dd\phi)_+\in\bbQ[v][p]$ be the quotient in the division $p\,\phi'(p)=(\cX\dd\phi)_+\,(p^5-v)+r$, $\deg r<5$; for $\phi=p^m$, $(\cX\dd\phi)_+=m\sum_{j\geq 0,\ m-5-5j\geq 0}v^jp^{m-5-5j}$. Then $(2\pi\ri)^{-1}\oint_{|p|=R}\phi\,\mathrm{Lie}_{\cX}w$ equals $-\langle(\cX\dd\phi)_+,w\rangle$.
\item $\int_{\mathfrak r}\mathrm{Lie}_{\cX}w=0$.
\end{enumerate}
\end{lem}

\begin{proof}
(i) On a curve $\mathrm{Lie}_{\cX}(\phi w)=\dd(\iota_{\cX}(\phi w))$, so $\oint\mathrm{Lie}_{\cX}(\phi w)=0$, and $\mathrm{Lie}_{\cX}(\phi w)=(\cX\phi)\,w+\phi\,\mathrm{Lie}_{\cX}w$. Here $\cX\phi=p\phi'/(p^5-v)=(\cX\dd\phi)_++r/(p^5-v)$, and $\oint\frac{r}{p^5-v}w=0$, since this form has all its poles inside the circle and is $O(\dd p/p^3)$ at $\infty$. The circle pairing of the polynomial $(\cX\dd\phi)_+$ is its residue pairing (\cref{def:mod-pair}). For $\phi=p^m$, $mp^m/(p^5-v)=m\sum_{j\geq 0}v^jp^{m-5-5j}$ for $|p|^5>|v|$, and the quotient is the polynomial part. (ii) $\int_{\mathfrak r}\mathrm{Lie}_{\cX}w=[\iota_{\cX}w]_0^\infty=0$: $\cX$ vanishes at $0$, where $w$ is holomorphic, and at $\infty$, $\cX=O(p^{-4})\,\dd/\dd p$ while $w=O(\dd p/p^2)$.
\end{proof}

\begin{thm}[The variational formula, in the form used]
\label{thm:mod-var}
Let $J=(p_1,\dots,p_n)$ be points off the branch points.
\begin{enumerate}
\item For $n\geq 1$ and $h+n>1$: $\partial_v|_x\,\omega_{h,n}(J)=\int_{q\in\mathfrak r}\omega_{h,n+1}(J,q)$.
\item For $h\geq 2$: $\partial_v\FTR_h=\int_{\mathfrak r}\omega_{h,1}$.
\end{enumerate}
\end{thm}

\begin{proof}
This is \cite[Thm.~5.1]{Eynard:2007kz}, (5-14) for $g+k>1$ and (5-15) for $g\geq 1$, for the temperature variation of \cite[\S5.3, (5-27)--(5-30)]{Eynard:2007kz}, translated by \cref{conv:mod-kernel}. In the normalisation of \cite{Eynard:2007kz} a variation of the curve at fixed $x$ is described by $\delta_\Omega y|_x\,\dd x=-\Omega$ \cite[(5-1)]{Eynard:2007kz} with $\Omega(p)=\int_{\partial\Omega}B(p,q)\Lambda(q)$ \cite[(5-13)]{Eynard:2007kz}, and (5-14), (5-15) read $\delta_\Omega W^{(h)}_n(J)|_{x}=\int_{\partial\Omega}W^{(h)}_{n+1}(J,q)\Lambda(q)$ and $\delta_\Omega F^{(h)}=-\int_{\partial\Omega}W^{(h)}_1\Lambda$. By \cref{lem:mod-fixedx}, $\Omega=-\partial_v(y\,\dd x)|_x=\dd p/p=-\int_{\mathfrak r}\omega_{0,2}(\cdot,q)$: the path $\partial\Omega=\mathfrak r$ joins the two poles $0$ and $\infty$ of $y\,\dd x$, and $\Lambda=-1$ on it. With $\omega_{h,n}=(-1)^nW^{(h)}_n$ and $\FTR_h=-F^{(h)}$ this gives $\partial_v|_x\omega_{h,n}=(-1)^n\int_{\mathfrak r}W^{(h)}_{n+1}\Lambda=(-1)^{n+1}(-1)^{n+1}\int_{\mathfrak r}\omega_{h,n+1}=\int_{\mathfrak r}\omega_{h,n+1}$ and $\partial_v\FTR_h=-\delta_\Omega F^{(h)}=\int_{\mathfrak r}W^{(h)}_1\Lambda=-\int_{\mathfrak r}W^{(h)}_1=\int_{\mathfrak r}\omega_{h,1}$. The hypotheses hold: the branch points are simple with $\dd y\neq 0$ for every $v\neq 0$; the variation is at fixed $x$ and the Bergman kernel is that of the fixed $\bbP^1$; $\Omega/\dd y=1/p$ is regular at the branch points, as required for \cite[(5-4)]{Eynard:2007kz}; the path avoids the branch points and has no hard edges \cite[(5-13), footnote~8]{Eynard:2007kz}; the integrals converge, since $\omega_{h,n+1}$ is holomorphic at $0$ and $\infty$ in each variable; and the residues at the moving branch points commute with $\partial_v|_x$, computed on fixed small circles in the $x$-plane around the moving critical values. The case $(h,n)=(0,1)$, which (i) excludes, is the hypothesis $\partial_v(y\,\dd x)|_x=\int_{\mathfrak r}\omega_{0,2}$ itself.
\end{proof}

The case $(h,n)=(0,2)$ of (i) is also proved directly, in \cref{lem:mod-genus0}~(iii) below. Part~(ii) is used in no proof.

\begin{lem}[Genus zero]
\label{lem:mod-genus0}
Put $e_{b,k}(p):=\dd p/(p-\varsigma_b)^{k+1}$ and $r_b:=1/(5\varsigma_b^3)$. Then:
\begin{enumerate}
\item $\partial_v|_x\,\omega_{0,2}(p_1,p_2)=\sum_br_b\,e_{b,1}(p_1)\,e_{b,1}(p_2)$;
\item $\omega_{0,3}(p_0,p_1,p_2)=-\sum_be_{b,1}(p_0)e_{b,1}(p_1)e_{b,1}(p_2)/x''(\varsigma_b)$, the negative of \cite[Thm.~4.1 (4-20)]{Eynard:2007kz};
\item $\int_{q\in\mathfrak r}\omega_{0,3}(p_1,p_2,q)=\partial_v|_x\,\omega_{0,2}(p_1,p_2)$;
\item for $\phi_1,\phi_2,\phi_3\in\bbQ[v][p]$,
\[
\begin{aligned}
\langle\phi_1\phi_2\phi_3\rangle^{\mathrm{TR}}_0&=-\tfrac15\textstyle\sum_b\varsigma_b^{-2}\prod_i\phi_i'(\varsigma_b), &\langle\phi_1\phi_2,-\log p\rangle^{\mathrm{TR}}_0&=\tfrac15\textstyle\sum_b\varsigma_b^{-3}\phi_1'(\varsigma_b)\phi_2'(\varsigma_b),\\
\langle\phi_1,(-\log p)^2\rangle^{\mathrm{TR}}_0&=-\tfrac15\textstyle\sum_b\varsigma_b^{-4}\phi_1'(\varsigma_b), &\langle(-\log p)^3\rangle^{\mathrm{TR}}_0&=1/v,
\end{aligned}
\]
and every pairing of $\omega_{0,3}$ with at least one leg in $\bbQ[v][p]$ lies in $\bbQ[v]$.
\end{enumerate}
\end{lem}

\begin{proof}
(i) $\omega_{0,2}$ does not depend on $v$ at fixed $p$, so by \cref{lem:mod-fixedx}, $\partial_v|_x\omega_{0,2}=-(\mathrm{Lie}_{\cX,1}+\mathrm{Lie}_{\cX,2})\,\omega_{0,2}=-\dd_1\dd_2\big[(\cX(p_1)-\cX(p_2))/(p_1-p_2)\big]$, with $\cX$ read as the function $p/(p^5-v)$; indeed, for a function $g$ and the field $g\,\dd/\dd p$ acting in both variables,
\[
(\mathrm{Lie}_{1}+\mathrm{Lie}_{2})\frac{\dd p_1\dd p_2}{(p_1-p_2)^2}=\Big(\frac{g'(p_1)+g'(p_2)}{(p_1-p_2)^2}-\frac{2(g(p_1)-g(p_2))}{(p_1-p_2)^3}\Big)\dd p_1\dd p_2=\dd_1\dd_2\frac{g(p_1)-g(p_2)}{p_1-p_2}.
\]
By partial fractions $p/(p^5-v)=\sum_br_b/(p-\varsigma_b)$, with no polynomial part, so $(\cX(p_1)-\cX(p_2))/(p_1-p_2)=-\sum_br_b/((p_1-\varsigma_b)(p_2-\varsigma_b))$, whence (i).
(ii) Only the two terms of ${\sum}'$ with $(h_1,I_1)=(0,\{p_1\})$, $(0,\{p_2\})$ occur in \eqref{eq:mod-kernel}. Near $\varsigma_b$ let $t:=q-\varsigma_b$, so $\iota_bq=\varsigma_b-t+O(t^2)$ and $\dd x(q)=x''(\varsigma_b)\,t(1+O(t))\,\dd t$. Then $\int_{\iota_bq}^q\omega_{0,2}(p_0,\cdot)=\dd p_0\Big(\frac1{p_0-q}-\frac1{p_0-\iota_bq}\Big)=\frac{2t\,\dd p_0}{(p_0-\varsigma_b)^2}(1+O(t))$ and $y(q)-y(\iota_bq)=2t(1+O(t))$, so $K_b(p_0,q)=\frac{e_{b,1}(p_0)}{2x''(\varsigma_b)t\,\dd t}(1+O(t))$. The bracket is $\omega_{0,2}(q,p_1)\omega_{0,2}(\iota_bq,p_2)+\omega_{0,2}(\iota_bq,p_1)\omega_{0,2}(q,p_2)=-2\,e_{b,1}(p_1)e_{b,1}(p_2)\,\dd t^2(1+O(t))$, because $\dd(\iota_bq)=-\dd t(1+O(t))$. The residue at $t=0$ of the product is $-e_{b,1}(p_0)e_{b,1}(p_1)e_{b,1}(p_2)/x''(\varsigma_b)$.
(iii) Since $x''(\varsigma_b)=5\varsigma_b^2$ and
\[
\int_{\mathfrak r}e_{b,1}=\Big[-\frac{1}{p-\varsigma_b}\Big]_0^\infty=-\frac{1}{\varsigma_b},
\]
(ii) gives $\int_{q\in\mathfrak r}\omega_{0,3}(p_1,p_2,q)=\sum_br_b\,e_{b,1}(p_1)\,e_{b,1}(p_2)$, which is (i).
(iv) $\Res_{p=\varsigma_b}\phi(p)\,e_{b,1}(p)=\phi'(\varsigma_b)$ and $\langle-\log p,e_{b,1}\rangle=\int_{\mathfrak r}e_{b,1}=-1/\varsigma_b$; inserting into (ii) gives the four formulas, the last one being $\frac15\sum_b\varsigma_b^{-5}=1/v$. Constants pair to zero, so in a pairing with an L-leg we may assume $\phi_i\in p\,\bbQ[v][p]$; then the powers of $\varsigma_b$ that occur are $\geq -4$, and $\sum_b\varsigma_b^N=5v^{N/5}$ if $5\mid N$ and $0$ otherwise. So only non-negative powers of $v$ survive.
\end{proof}

\begin{prop}[The recursion (V1)]
\label{prop:mod-V1}
Let $h\geq 0$, $a,k\geq 0$, $n=a+k\geq 1$, $2h-2+n>0$, and $\phi_1,\dots,\phi_a\in\bbQ[v][p]$. Then
\begin{equation}
\label{eq:mod-V1}
\begin{aligned}
\langle\phi_1\cdots\phi_a,(-\log p)^{k+1}\rangle^{\mathrm{TR}}_h={}&\partial_v\langle\phi_1\cdots\phi_a,(-\log p)^k\rangle^{\mathrm{TR}}_h-\sum_{i=1}^a\langle\cdots,\partial_v\phi_i,\cdots,(-\log p)^k\rangle^{\mathrm{TR}}_h\\
&+\sum_{i=1}^a\langle\cdots,(\cX\dd\phi_i)_+,\cdots,(-\log p)^k\rangle^{\mathrm{TR}}_h ,
\end{aligned}
\end{equation}
where $\partial_v\phi_i$ is the derivative of the coefficients and the dots stand for the other legs.
\end{prop}

\begin{proof}
Fix $v_0\neq 0$. For $v$ near $v_0$ choose $R>\max_b|\varsigma_b(v)|$ and a ray $\mathfrak r$ avoiding all $\varsigma_b(v)$; then $C(v):=\langle\phi_1\cdots\phi_a,(-\log p)^k\rangle^{\mathrm{TR}}_h=(2\pi\ri)^{-a}\oint_{|p_i|=R}\int_{\mathfrak r^k}\prod_i\phi_i(p_i)\,\omega_{h,n}$ with contours independent of $v$. The integrand is holomorphic in $v$ on the compact contours and, on the rays, bounded near $v_0$ uniformly in $v$ by a constant times $|\dd q|/(1+|q|)^2$ (holomorphy at $0$, decay at $\infty$, continuity of the coefficients). So we may differentiate under the integrals: $\partial_vC=\sum_i\langle\cdots,\partial_v\phi_i,\cdots\rangle^{\mathrm{TR}}_h+$ the integral of $\prod_i\phi_i\,\partial_v|_p\omega_{h,n}$. By \cref{lem:mod-fixedx} and \cref{thm:mod-var}~(i) (applicable since $n\geq 1$ and $h+n>1$), $\partial_v|_p\omega_{h,n}=\int_{q\in\mathfrak r}\omega_{h,n+1}(\cdot,q)+\sum_{l=1}^n\mathrm{Lie}_{\cX,l}\,\omega_{h,n}$. The first term gives the pairing with one more log leg (Fubini; there are no diagonal poles); a Lie term in an L-slot gives $-\langle\cdots,(\cX\dd\phi_i)_+,\cdots\rangle^{\mathrm{TR}}_h$ by \cref{lem:mod-ibp}~(i), and in a log slot it gives $0$ by \cref{lem:mod-ibp}~(ii). Rearranging gives \eqref{eq:mod-V1}.
\end{proof}

% Section 6, part 9.
The base case $h=1$ of Theorem~V rests on a local formula for $\omega_{1,1}$. Let $c$ be a branch point, $t:=p-c$, and $\chi$ (a local letter) a local coordinate at $c$ with $x=x(c)+\lambda\chi^2$, $\lambda\neq 0$ a constant, so that $\iota_c\colon\chi\mapsto-\chi$. Write $y=\sum_ky_k\chi^k$ and, near $(c,c)$, $\omega_{0,2}=\big[(\chi_1-\chi_2)^{-2}+\mathrm{reg}(\chi_1,\chi_2)\big]\dd\chi_1\dd\chi_2$ with $\mathrm{reg}$ holomorphic and symmetric; then $\mathrm{reg}_{00}:=\mathrm{reg}(0,0)=\frac16\{t;\chi\}(0)$, where $\{t;\chi\}=t'''/t'-\frac32(t''/t')^2$ is the Schwarzian derivative of $t$ as a function of $\chi$ (the regular part on the diagonal of $t'(\chi_1)t'(\chi_2)/(t(\chi_1)-t(\chi_2))^2$). For a $1$-form $\eta=\sum_m\eta_m\chi^m\dd\chi$ holomorphic at $c$ put
\begin{equation}
\label{eq:mod-loc}
\mathrm{Loc}_c[\eta]:=-\frac{1}{4\lambda y_1}\Big[\mathrm{reg}_{00}\,\eta_0+\frac{\eta_2}{12}-\frac{y_3}{4y_1}\,\eta_0\Big].
\end{equation}

\begin{lem}[A local formula for $\omega_{1,1}$]
\label{lem:mod-w11}
\begin{enumerate}
\item $\omega_{1,1}(p_0)=\sum_c\mathrm{Loc}_c[\omega_{0,2}(p_0,\cdot)]$ for $p_0$ off the branch points, $\omega_{0,2}(p_0,\cdot)$ being a form in the second slot.
\item $\omega_{1,1}=\dd G$ with $G$ rational, with poles only at the branch points and $G(\infty)=0$; so $\omega_{1,1}$ has zero residues, is holomorphic at $p=0$ and is $O(\dd p/p^2)$ at $\infty$.
\item $\langle-\log p\rangle^{\mathrm{TR}}_1=\int_{\mathfrak r}\omega_{1,1}=\sum_c\mathrm{Loc}_c[-\dd q/q]$.
\end{enumerate}
\end{lem}

\begin{proof}
(i) By \eqref{eq:mod-kernel},
\[
\omega_{1,1}(p_0)=\sum_c\Res_{q\to c}K_c(p_0,q)\,\omega_{0,2}(q,\iota_cq).
\]
Expand $\omega_{0,2}(p_0,q')=\dd p_0\sum_mB_m(p_0)\chi^m\dd\chi$ for $q'$ near $c$ with coordinate $\chi$. Then $\int_{\iota_cq}^q\omega_{0,2}(p_0,\cdot)=2\,\dd p_0\sum_{m\text{ even}}B_m\chi^{m+1}/(m+1)$, $y(q)-y(\iota_cq)=2\sum_{k\text{ odd}}y_k\chi^k$ and $\dd x=2\lambda\chi\,\dd\chi$, so
\[
K_c(p_0,q)=\frac{\dd p_0}{\dd\chi}\,\frac{1}{4\lambda y_1}\Big[B_0\chi^{-1}+\Big(\frac{B_2}{3}-B_0\frac{y_3}{y_1}\Big)\chi+O(\chi^3)\Big],
\]
and $\omega_{0,2}(q,\iota_cq)=-\big[(4\chi^2)^{-1}+\mathrm{reg}_{00}+O(\chi^2)\big]\dd\chi^2$, since $\dd(-\chi)=-\dd\chi$ and $\mathrm{reg}(\chi,-\chi)=\mathrm{reg}_{00}+O(\chi^2)$ by symmetry. The coefficient of $\chi^{-1}\dd\chi$ in the product is
\[
-\frac{1}{4\lambda y_1}\Big[B_0\,\mathrm{reg}_{00}+\frac14\Big(\frac{B_2}{3}-B_0\frac{y_3}{y_1}\Big)\Big]\dd p_0=\mathrm{Loc}_c[\omega_{0,2}(p_0,\cdot)].
\]
(ii) $\omega_{0,2}(p_0,q)=\dd_{p_0}\big[\dd q/(q-p_0)\big]$, and near $c$, $\dd q/(q-p_0)=\sum_mg_m(p_0)\chi^m\dd\chi$ with $g_m$ rational in $p_0$, with poles only at $c$ and $g_m(\infty)=0$ (expand $1/(c+t(\chi)-p_0)$ in powers of $t$). So $B_m(p_0)\dd p_0=\dd g_m(p_0)$, and by (i) $\omega_{1,1}$ is the differential of a linear combination $G$ of the $g_0,g_2$ at the five branch points. (iii) Each residue in (i) is an integral over a small circle around $c$, which $\mathfrak r$ avoids; by (ii) the ray integral converges, and since $\omega_{0,2}(p_0,q)=O(p_0^{-2})\dd p_0$ uniformly for $q$ on the circles, the integral over $p_0\in\mathfrak r$ may be taken inside: $\int_{p_0\in\mathfrak r}\omega_{0,2}(p_0,q)=\dd q\big[-1/(p_0-q)\big]_{p_0=0}^{p_0=\infty}=-\dd q/q$ for $q$ off the ray, and $\mathrm{Loc}_c$ is linear in finitely many Taylor coefficients.
\end{proof}

By (i) and \eqref{eq:mod-kernel}, $\mathrm{Loc}_c[\eta]=\Res_{q\to c}\frac12\frac{\int_{\iota_cq}^q\eta}{(y(q)-y(\iota_cq))\dd x(q)}\,\omega_{0,2}(q,\iota_cq)$ for $\eta$ holomorphic at $c$; this form does not depend on the coordinate $\chi$.

\begin{thm}[Theorem V]
\label{thm:mod-V}
Let $h\geq 0$, $a,k\geq 0$, $n=a+k\geq 1$, $2h-2+n>0$, and $\phi_1,\dots,\phi_a\in\bbQ[v][p]$.
\begin{enumerate}
\item If $a\geq 1$, then $\langle\phi_1\cdots\phi_a,(-\log p)^k\rangle^{\mathrm{TR}}_h\in\bbQ[v]$.
\item If $a=0$, then $\langle(-\log p)^k\rangle^{\mathrm{TR}}_h=(\FTR_h)^{(k)}(v)$, with $\FTR_h$ of \cref{thm:mod-gap} for $h\geq 2$, $\FTR_1:=-\frac1{12}\log v$ (here $k\geq 1$) and $\FTR_0:=\frac12v^2\log v$ (here $k\geq 3$).
\item The recursion \eqref{eq:mod-V1} holds.
\end{enumerate}
\end{thm}

The quoted inputs of the proof are \cite[Thm.~5.1, (5-14)]{Eynard:2007kz} (through \cref{thm:mod-var}~(i) and \cref{prop:mod-V1}) and \cite[Thm.~7.1]{ABDKS25} (through \cref{thm:mod-Lreg} and \cref{thm:mod-gap}), besides \cite[Thm.~4.2, Thm.~4.6]{Eynard:2007kz} (through the paragraph after \cref{conv:mod-kernel}, \cref{def:mod-pair} and \cref{lem:mod-xfun}).

\begin{proof}
(iii) is \cref{prop:mod-V1}. (i) Fix $h$ and $a\geq 1$; we use induction on $k$. If $2h-2+a>0$, the case $k=0$ is \cref{thm:mod-Lreg}. If $2h-2+a\leq 0$, then $(h,a)$ is $(0,1)$ or $(0,2)$, the first admissible pairing is $\langle\phi_1,(-\log p)^2\rangle^{\mathrm{TR}}_0$, respectively $\langle\phi_1\phi_2,-\log p\rangle^{\mathrm{TR}}_0$, and it lies in $\bbQ[v]$ by \cref{lem:mod-genus0}~(iv). For the step from $k$ to $k+1$ (with $2h-2+a+k>0$), \eqref{eq:mod-V1} writes the pairing with $k+1$ log legs through the $v$-derivative of a pairing with $k$ log legs and pairings with the same number $a$ of L-legs and $k$ log legs, the legs $\partial_v\phi_i$ and $(\cX\dd\phi_i)_+$ lying in $\bbQ[v][p]$ because $p^5-v$ is monic; all of these lie in $\bbQ[v]$ by the induction hypothesis. The L-legs are never consumed, so the induction stays in the class $a\geq 1$.

(ii) For $a=0$, \eqref{eq:mod-V1} reads $\langle(-\log p)^{k+1}\rangle^{\mathrm{TR}}_h=\partial_v\langle(-\log p)^k\rangle^{\mathrm{TR}}_h$, so it suffices to treat the first admissible $k$. \emph{$h\geq 2$, $k=1$.} As in the proof of \cref{prop:mod-period}~(iii), $(2-2h)\FTR_h=\frac15\langle p^5\rangle^{\mathrm{TR}}_h+v\langle-\log p\rangle^{\mathrm{TR}}_h$ with $\langle p^5\rangle^{\mathrm{TR}}_h=0$; since $\FTR_h$ is a multiple of $v^{2-2h}$ (\cref{thm:mod-gap}), $\langle-\log p\rangle^{\mathrm{TR}}_h=(2-2h)\FTR_h/v=(\FTR_h)'(v)$. This case uses neither \cref{thm:mod-var} nor anything beyond \cref{thm:mod-Lreg} and \cref{thm:mod-gap}. \emph{$h=0$, $k=3$.} $\langle(-\log p)^3\rangle^{\mathrm{TR}}_0=1/v=(\FTR_0)'''$ by \cref{lem:mod-genus0}~(iv). \emph{$h=1$, $k=1$.} We show $\langle-\log p\rangle^{\mathrm{TR}}_1=-1/(12v)$.

\emph{Reduction to $v=1$.} By \cref{lem:mod-scale}~(ii) with $k=\varsigma$, $m_\varsigma^*\omega^{(v)}_{1,1}=\varsigma^{5(2-2-1)}\omega^{(1)}_{1,1}=v^{-1}\omega^{(1)}_{1,1}$; $m_\varsigma$ maps rays to rays, and the pairing does not depend on the ray. So $\langle-\log p\rangle^{\mathrm{TR}}_1(v)=v^{-1}\langle-\log p\rangle^{\mathrm{TR}}_1(1)$. Let $v=1$, $\varsigma=1$, so $\varsigma_b=\zeta^{-b}$.

\emph{Five equal contributions.} The map $\rho_b(p):=\zeta^{-b}p$ sends $1$ to $\varsigma_b$ and satisfies $x\circ\rho_b=\zeta^bx$, $y\circ\rho_b=\zeta^{-b}y$, $\rho_b^*\omega_{0,2}=\omega_{0,2}$ and $\rho_b^*(\dd q/q)=\dd q/q$; hence $\iota_b\circ\rho_b=\rho_b\circ\iota_0$, and the factor $(y(q)-y(\iota q))\,\dd x(q)$ is unchanged ($\zeta^{-b}\zeta^b=1$). By the coordinate-free form of $\mathrm{Loc}$ after \cref{lem:mod-w11}, $\mathrm{Loc}_{\varsigma_b}[-\dd q/q]=\mathrm{Loc}_1[-\dd q/q]$ for every $b$.

\emph{The contribution at $c=1$.} With $t=p-1$, $x(1+t)=\frac54+\frac52t^2+0\cdot t^3+\frac54t^4-t^5+O(t^6)$, so $\lambda=5/2$ and $\chi=t\big(1+\frac{t^2}2-\frac{2t^3}5+O(t^4)\big)^{1/2}=t+\frac{t^3}4-\frac{t^4}5+O(t^5)$, that is $t=\chi-\frac{\chi^3}4+\frac{\chi^4}5+O(\chi^5)$. Hence $y=1+t$ has $y_1=1$, $y_2=0$, $y_3=-\frac14$; $t'(0)=1$, $t''(0)=0$, $t'''(0)=-\frac32$, so $\{t;\chi\}(0)=-\frac32$ and $\mathrm{reg}_{00}=-\frac14$; and $-\dd q/q=-t'(\chi)\,\dd\chi/(1+t(\chi))=-\big(1-\chi+\frac14\chi^2+O(\chi^3)\big)\dd\chi$, so $\eta_0=-1$, $\eta_2=-\frac14$. By \eqref{eq:mod-loc},
\[
\mathrm{Loc}_1[-\dd q/q]=-\frac1{10}\Big[\Big(-\frac14\Big)(-1)+\frac{-1/4}{12}-\frac{-1/4}{4}\,(-1)\Big]=-\frac1{10}\cdot\frac{12-1-3}{48}=-\frac1{60}.
\]
By \cref{lem:mod-w11}~(iii) the total at $v=1$ is $5\cdot(-\frac1{60})=-\frac1{12}$. So $\langle-\log p\rangle^{\mathrm{TR}}_1=-1/(12v)=(\FTR_1)'(v)$, which also fixes $\FTR_1$ up to an additive constant.
\end{proof}

\input{sections/s06_model_10}



